Con il presente lavoro di tesi sono stati ottenuti dei risultati, che verranno delineati nella Sezione 6.1, esaminandoli criticamente e mettendo in evidenza le realizzazioni e le scoperte emerse durante l'analisi, lo sviluppo e la valutazione delle soluzioni proposte. La Sezione 6.2 fornirà uno sguardo approfondito sui limiti di questo lavoro e sulle sfide incontrate durante il corso di questo progetto. La Sezione 6.3 analizzerà le prospettive future del presente lavoro, suggerendo vie potenziali per ulteriori sviluppi. Questa Sezione non solo proietterà lo sguardo oltre la conclusione di questo progetto, ma fornirà anche una base per ulteriori contributi e iterazioni sulle soluzioni proposte.

\section{Risultati ottenuti}
Grazie al lavoro di analisi, progettazione, sviluppo e test, questa tesi ha portato alla creazione di una nuova versione dell'applicazione ZConnect, migliorata secondo le richieste dell'Azienda X nella sua funzionalità di monitoraggio delle presenze. La collaborazione con il team di sviluppo di ZConnect e il team di test ha permesso il rilascio della versione migliorata, attualmente disponibile sugli store di Android ed iOS, con un ampio utilizzo da parte di centinaia di migliaia di utenti.

La recente evoluzione della funzionalità di monitoraggio delle presenze è stata guidata dalle esigenze specifiche dell'Azienda X, la quale ha espresso piena soddisfazione per il lavoro svolto. I miglioramenti apportati riguardano principalmente la mappa di ZConnect, ora progettata in modo da rendere chiare all'utente sia la propria posizione che quella dell'area di lavoro più vicina. Per fornire un maggior supporto all'utente nel raggiungere l'area di lavoro, è stata migliorata l'usabilità della mappa e introdotta la funzionalità "follow-me", ideata per guidare l'utente in modo più intuitivo verso il luogo di lavoro.

\section{Limiti del sistema}
L'analisi dei risultati di questo lavoro ha rivelato alcune limitazioni.
In primo luogo, la mancanza di test di usabilità con gli utenti rappresenta una limitazione significativa. Questi test avrebbero fornito un feedback diretto sull'efficacia delle modifiche apportate al sistema di monitoraggio delle timbrature, valutando se le soluzioni implementate hanno effettivamente risolto i problemi dei dipendenti dell'Azienda X, senza introdurne di nuovi.

Una seconda limitazione è emersa dalla mancanza di test reali di utilizzo da parte del team di test. La verifica statica su dispositivi fisici non ha coperto l'aspetto dinamico dello spostamento dell'utente sulla mappa, trascurando l'efficacia della funzionalità "follow-me". La correttezza e l'utilità di questa caratteristica non sono state completamente validate, rappresentando un'area di incertezza.

Infine, la struttura stessa di ZConnect rappresenta un limite, poiché, come già visto, è sviluppata sia per Android che per iOS, duplicando le basi di codice. I problemi affrontati, e risolti, da questo lavoro di tesi avrebbero potuto essere gestiti più efficientemente da un'applicazione ibrida. Questo avrebbe potenzialmente ridotto i tempi di sviluppo e semplificato la gestione complessiva del software, in quanto un'app ibrida consente di scrivere le funzionalità dell'app in un'unica base di codice \cite{AppNative}.  

\section{Miglioramenti futuri}
L'analisi dei risultati di questo lavoro di tesi ha evidenziato potenziali miglioramenti della funzionalità di monitoraggio delle presenze di ZConnect.

In primo luogo, si potrebbe implementare l'inserimento di etichette esplicative al momento della pressione del marker rappresentante il centro del luogo di lavoro. Questa caratteristica aggiuntiva consentirebbe agli utenti di ottenere informazioni chiare e dettagliate sul luogo di lavoro selezionato.

In secondo luogo, considerando che un utente potrebbe avere numerose aree di lavoro associate, sarebbe utile consentire la ricerca di un luogo di lavoro specifico in base al suo nome. Tale funzionalità permetterebbe all'utente di posizionare la mappa direttamente sul luogo di lavoro desiderato, semplificando l'utilizzo della mappa.

Infine, un ulteriore miglioramento potrebbe consistere nell'introduzione di un sistema di navigazione. Una volta selezionato un luogo di lavoro, questa funzione consentirebbe all'utente di visualizzare sulla mappa il percorso ottimale per raggiungere la destinazione desiderata. Questo sistema di navigazione potrebbe notevolmente facilitare gli spostamenti degli utenti verso le varie sedi di lavoro associate, migliorando complessivamente l'esperienza di utilizzo della funzionalità di monitoraggio delle timbrature di ZConnect.